\documentclass[a4paper,12pt]{article}
\usepackage{color}
\usepackage{graphicx, gensymb}
\usepackage{amsfonts, cancel, amssymb}
\usepackage{amsmath, amsthm, array, esvect, esint}
\usepackage{typearea}
\usepackage[a4paper,left=2cm,right=2cm,top=2cm,bottom=3cm,bindingoffset=5mm]{geometry}
\newcommand\numberthis{\addtocounter{equation}{1}\tag{\theequation}}
\newcommand{\bra}{}
\newcommand{\kom}{}
\newcommand{\brak}{}
\newcommand{\braket}{}
\newcommand{\intx}{}
\newcommand{\intr}{}
\newcommand{\kla}{}
\newcommand{\klam}{}
\newcommand{\klammer}{}
\newcommand{\oper}{}
\newcommand{\tecxt}{}
\newcommand{\Bra}{}
\newcommand{\Ket}{}

\begin{document}

\title{8 Ferromagnetismus}
\author{Joachim Schwardt}
\date{\today}
\maketitle

\section*{Aufgabe 8.1: Ferromagnetismus, Curie-Weiss-Modell}
Gegeben seien $N$ unterscheidbare Spins $\sigma_k\in\{\pm 1\}$, die alle untereinander wechselwirken. Ein externes Magnetfeld sei durch den Parameter $b$ beschrieben, die Wechselwirkung der Spins durch $J$. Also lautet der Hamiltonian:
\begin{align*}
H&=-\frac{J}{2N}\sum_{j,k}\sigma_j\sigma_k - b\sum_k\sigma_k
\end{align*}
Für eine mittlere Magnetisierung $\bra\langle {m}\rangle=\frac{1}{N}\bra\langle {M}\rangle$ mit $M=\sum_k\sigma_k$ gilt dann:
\begin{align*}
H&=-\frac{J}{2N}\sum_j\sigma_j\sum_k\sigma_k-b\sum_k\sigma_k=-\frac{J}{2N}\bra\langle {M}\rangle^2-b\bra\langle {M}\rangle=-N\kla\left( {\frac{J}{2}\bra\langle {m}\rangle+b\bra\langle {m}\rangle} \right)\propto N
\end{align*}
Es gibt $|\Omega|=2^N$ mögliche Mikrozustände. Die Anzahl der unterschiedlichen Energien hängt nur davon ab, wie viele Spins insgesamt positiv oder negativ sind. Dafür gibt es nur $N+1$ Möglichkeiten. Betrachte folgende Umformung:
\begin{align*}
2^N&=\kla\left( {1+1} \right)^N=\sum_{k=0}^N\frac{N!}{k!(N-k)!}
\end{align*}
Man kann die gesamte Anzahl der Mikrozustände also in $N+1$ Energiezustände zerlegen. Sei $k$ die Anzahl an Spins mit Wert $\sigma_k=+1$. Dann gibt es $\binom{N}{k}$ mögliche Mikrozustände.\\
Schreibe $B:=\beta J\sum_j\sigma_j$ und $A:=\frac{N\beta J}{2}$. Dann gilt für die Zustandssumme mit $e^{B^2/4A}=\sqrt{\frac{A}{\pi}}\intr\int_{\mathbb{R}^{ }}{e^{-Ax^2+Bx}}\,\mathrm{d}^{ }x$ im kanonischen Ensemble:
\begin{align*}
Z&=\sum_{\klammer\left\{{\sigma} \right\}} e^{-\beta H}=\sum_{\klammer\left\{{\sigma} \right\}}\exp\kla\left( {\frac{B^2}{2N\beta J}+\beta b\sum_k\sigma_k} \right)=\\
&=\sum_{\klammer\left\{{\sigma} \right\}}\sqrt{\frac{N\beta J}{2\pi}}\intr\int_{\mathbb{R}^{ }}{\exp\kla\left( {-\frac{N\beta J}{2}x^2+\beta(Jx+b)\sum_j\sigma_j} \right)}\,\mathrm{d}^{ }x=\\
&=\sqrt{\frac{N\beta J}{2\pi}}\intr\int_{\mathbb{R}^{ }}{\exp\kla\left( {-\frac{N\beta J}{2}x^2} \right)\sum_{\sigma_1\in\klammer\left\{{\pm 1} \right\}}\dots\sum_{\sigma_N\in\klammer\left\{{\pm 1} \right\}}\prod_k\exp\kla\left( {\beta(Jx+b)\sigma_k} \right)}\,\mathrm{d}^{ }x=\\
&=\sqrt{\frac{N\beta J}{2\pi}}\intr\int_{\mathbb{R}^{ }}{\exp\kla\left( {-\frac{N\beta J}{2}x^2} \right)\klam\left[ {2\cosh\kla\left( {\beta (Jx+b)} \right)} \right]^N}\,\mathrm{d}^{ }x
\end{align*}
Mit $L(x):=\frac{J}{2}x^2-\frac{1}{\beta}\log\klam\left[ {2\cosh(\beta(Jx+b))} \right]$ kann dieses Ergebnis umgeschrieben werden:
\begin{align*}
Z&=\sqrt{\frac{N\beta J}{2\pi}}\intr\int_{\mathbb{R}^{ }}{\exp\kla\left( {-N\beta L(x)} \right)}\,\mathrm{d}^{ }x
\end{align*}
Das Minimum $x_0$ der Funktion $L(x)$ ist durch eine implizite Gleichung gegeben:
\begin{align*}
0&\overset{!}{=}L'(x_0)&\Rightarrow &&Jx_0&=\frac{1}{\beta}\frac{\beta J\sinh(\beta(Jx_0+b))}{\cosh(\beta(Jx_0+b))}\\
\Rightarrow\ \ \ x_0&=\tanh\beta (Jx_0+b)
\end{align*}
Das Taylor-Polynome lautet in zweiter Ordnung:
\begin{align*}
L(x)&=L(x_0)+\frac{1}{2}(x-x_0)^2\klam\left[ {J-J\tanh'(\beta(Jx+b))\big\vert_{x=x_0}} \right]+\mathcal{O}(x^3)\approx\\
&\approx L(x_0)+\frac{J}{2}(x-x_0)^2\klam\left[ {1-\beta J\kla\left( {1-\tanh^2\beta(Jx_0+b)} \right)} \right]=\\
&=L(x_0)+\frac{J}{2}(x-x_0)^2\klam\left[ {1-\beta J+\beta J x_0^2} \right]
\end{align*}
Für große $N$ kann das Integral in der Zustandssumme dann wie folgt genähert werden:
\begin{align*}
Z&\approx\sqrt{\frac{N\beta J}{2\pi}}\exp\kla\left( {-N\beta L(x_0)} \right)\intr\int_{\mathbb{R}^{ }}{\exp\kla\left( {-\frac{J}{2}\klam\left[ {1-\beta J+\beta Jx_0^2} \right]x^2} \right)}\,\mathrm{d}^{ }x=\\
&=\sqrt{\frac{N\beta}{1-\beta J+\beta Jx_0^2}}\exp\kla\left( {-N\beta L(x_0)} \right)
\end{align*}
Betrachte nun den Grenzfall $b\rightarrow 0$, also ein verschwindendes äußeres Magnetfeld. Dann vereinfacht sich die implizite Gleichung zu $x_0=\tanh\beta Jx_0$. Eine (triviale) Lösung ist immer $x_0=0$. Falls der lineare Term um $x\simeq 0$ langsamer steigt als der hyperbolische, so gibt es noch zwei weitere Lösungen mit gleichem Betrag und unterschiedlichem Vorzeichen. Die Bedingung dafür ist $\beta J<1$ oder $J<k_BT$ mit kritischer Temperatur $T_c=\frac{J}{k_B}$.\\
Die freie Energie ist $F=-\frac{1}{\beta}\log Z$, woraus folgt:
\begin{align*}
f&:=\frac{F}{N}\approx L(x_0)-\frac{1}{2N\beta}\log\frac{N\beta}{1-\beta J+\beta Jx_0^2}=\\
&=\frac{J}{2}x_0^2-\frac{1}{\beta}\log\klam\left[ {2\cosh\beta (Jx_0+b)} \right]-\frac{1}{2N\beta}\log\frac{N\beta}{1-\beta J+\beta Jx_0^2}
\end{align*}
Im Grenzwert $N\rightarrow\infty $ gilt dann:
\begin{align*}
\lim f&=\frac{J}{2}x_0^2-\frac{\log 2}{\beta}+\frac{1}{2\beta}\log\klam\left[ {1-x_0^2} \right]
\end{align*}
Betrachte die implizite Gleichung unter Ableitung nach $b$:
\begin{align*}
\partial_bx_0&=\beta\klam\left[ {1+J\partial_bx_0} \right]\tecxt\text{sech}^2\beta(Jx_0+b)\\
\Rightarrow\ \ \ \partial_bx_0&=\frac{\beta\tecxt\text{sech}^2\beta(Jx_0+b)}{1-\beta J\tecxt\text{sech}^2\beta(Jx_0+b)}=\frac{\beta (1-x_0^2)}{1-\beta J(1-x_0^2)}
\end{align*}
Für den mittleren Spin gilt dann:
\begin{align*}
m&=\bra\langle {\sigma}\rangle=-\partial_b\lim f=-Jx_0\partial_bx_0+\frac{1}{\beta}\frac{x_0\partial_bx_0}{1-x_0^2}=\frac{x_0\partial_bx_0}{\beta}\frac{1-\beta J(1-x_0^2)}{1-x_0^2}=x_0
\end{align*}
Betrachte erneut die implizite Gleichung im Grenzwert $b\rightarrow 0$, diesmal mit $T_c:=\frac{J}{k_B}$ und $T=\frac{1}{k_B\beta}$. Entwickle in erster Ordnung von $(T-T_c)$ und dann in dritter Ordnung von $m$ mit $\tanh x\approx x-\frac{x^3}{3}$ und $\tecxt\text{sech}^2x\approx 1-x^2$:
\begin{align*}
m&=\tanh\frac{T_cm}{T}\approx\tanh m-(T-T_c)\frac{T_cm}{T^2}\tecxt\text{sech}^2m\approx\\
&\approx m-\frac{m^3}{3}-\kla\left( {1-\frac{T_c}{T}} \right)\frac{T_c}{T}m(1-m^2)
\end{align*}
Multipliziere beide Seiten mit $\frac{3T^2}{m}$:
\begin{align*}
T^2m^2&\approx 3\kla\left( {T_c-T} \right)T_c(1-m^2)\\
\Rightarrow\ \ \ m(T)&\approx\sqrt{\frac{3T_c(T_c-T)}{T^2+3T_c(T_c-T)}}=\sqrt{1-\frac{1}{1+3\frac{T_c}{T}\kla\left( {\frac{T_c}{T}-1} \right)}}
\end{align*}
\begin{align*}
m&=\tanh\frac{T_cm}{T}\approx\frac{T_cm}{T}-\frac{1}{3}\kla\left( {\frac{T_cm}{T}} \right)^3\\
\Rightarrow\ \ \ m(T)^2&\approx 3\kla\left( {\frac{T}{T_c}-\frac{T^2}{T_c^2}} \right)\approx \frac{3}{T_c}(T_c-T)\\
\Rightarrow\ \ \ m(T)&\approx m_0\kla\left( {T_c-T} \right)^{1/2}
\end{align*}
Dabei ist der kritische Exponent $c=\frac{1}{2}$ und der Parameter $m_0=\sqrt{\frac{3}{T_c}}$.\newpage
\section*{Aufgabe 8.2: Ferromagnetismus, 1D-Ising-Modell}
Der Hamiltonian ist bei periodischen Randbedingungen der Spins gegeben durch:
\begin{align*}
H&=-J\sum_k\sigma_k\sigma_{k+1}-b\sum_k\sigma_k
\end{align*}
Für die Zustandssumme gilt wie üblich:
\begin{align*}
Z&=\sum_{\klammer\left\{{\sigma} \right\}}\exp\kla\left( {-\beta H} \right)=\sum_{\klammer\left\{{\sigma} \right\}}\prod_k\exp\kla\left( {\beta J\sigma_k\sigma_{k+1}+\beta b\sigma_k} \right)=:\sum_{\{\sigma\}}\prod_kz_k
\end{align*}
Die vier möglichen Werte von $z_k$ können in einer Matrix dargestellt werden. Dabei zählen die Zeilen den Wert des Spins $\sigma_k$ durch und die Spalten den Wert von $\sigma_{k+1}$:
\begin{align*}
T&:=\begin{pmatrix}
e^{\beta(J+b)}&e^{\beta (-J+b)} \\ e^{\beta(-J-b)}&e^{\beta (J-b)}
\end{pmatrix}
\end{align*}
Die ``richtigen'' Spins wählt man durch Multiplikation mit Spin-Vektoren aus. Schreibe dazu $\vec{\sigma}_k=\binom{1}{0}$ für $\sigma_k=+1$ und $\vec{\sigma}_k=\binom{0}{1}$ für $\sigma_k=-1$. Dann gilt $z_k=\brak\langle {\vec{\sigma}_k}| {T|\vec{\sigma}_{k+1}}\rangle$. Wir können die Zustandssumme mit $\sum_{\{\sigma_j\}}\oper | {\vec{\sigma}_j} \rangle\langle {\vec{\sigma}_j} |=I$ dann wie folgt vereinfachen. Beachte noch, dass für allgemeine Matrizen $\tecxt\text{tr\,}M=\sum_{\{\sigma_k\}}\braket\langle {\vec{\sigma}_k} | {M} | {\vec{\sigma}_{k}}\rangle$ gilt:
\begin{align*}
Z&=\sum_{\{\sigma\}}\prod_k\brak\langle {\vec{\sigma}_k}| {T|\vec{\sigma}_{k+1}}\rangle=\sum_{\{\sigma_1\}}\braket\langle {\vec{\sigma}_1} | {T^N} | {\vec{\sigma}_{N+1}}\rangle=\tecxt\text{tr\,}T^N=\lambda_1^N+\lambda_2^N
\end{align*}
Wir benötigen also die Eigenwerte der Transfermatrix $T$:
\begin{align*}
0&\overset{!}{=}\det\kla\left( {T-\lambda I} \right)=\lambda^2-\klam\left[ {e^{\beta(J+b)}+e^{\beta(J-b)}} \right]\lambda +\klam\left[ {e^{2\beta J}-e^{-2\beta J}} \right]=\\
&=\lambda^2-2\lambda e^{\beta J}\cosh\beta b + e^{-2\beta J}-e^{-2\beta J}\\
\Rightarrow\ \ \ \lambda&=e^{\beta J}\cosh\beta b\pm\sqrt{e^{2\beta J}\cosh^2\beta b-e^{2\beta J}+e^{-2\beta J}}=\\
&=e^{\beta J}\klam\left[ {\cosh\beta b\pm\sqrt{\sinh^2\beta b+e^{-4\beta J}}} \right]
\end{align*}
Der Eigenwert mit positivem Vorzeichen ist dabei größer als der mit negativem. Im Grenzwert $N\rightarrow\infty $ folgt dann für die Zustandssumme:
\begin{align*}
Z&=\lambda_1^N\kla\left( {1+\frac{\lambda_2^N}{\lambda_1^N}} \right)=\lambda_1^N\klam\left[ {1+\kla\left( {\frac{\cosh\beta b -\sqrt{\sinh^2\beta b+e^{-4\beta J}}}{\cosh\beta b +\sqrt{\sinh^2\beta b+e^{-4\beta J}}}} \right)^N} \right]\longrightarrow\lambda_1^N
\end{align*}
Die freie Energie pro Spin ist dann gegeben durch:
\begin{align*}
f&=-\frac{1}{\beta N}\log Z\simeq -\frac{1}{\beta}\log\lambda_1=-J -\frac{1}{\beta}\log\klam\left[ {\cosh\beta b+\sqrt{\sinh^2\beta b+e^{-4\beta J}}} \right]
\end{align*}
Die mittlere Magnetisierung pro Spin folgt dann zu:
\begin{align*}
m&=\bra\langle {\sigma}\rangle=-\partial_bf=\frac{\sinh\beta b+\frac{\sinh\beta b\cosh\beta b}{\sqrt{\sinh^2\beta b+e^{-4\beta J}}}}{\cosh\beta b+\sqrt{\sinh^2\beta b+e^{-4\beta J}}}=\frac{\sinh\beta b}{\sqrt{\sinh^2\beta b+e^{-4\beta J}}}\kom\overset{\substack{{b\rightarrow 0} \\ |}}{\longrightarrow}0
\end{align*}
Im Grenzwert $b\rightarrow 0$ vereinfacht sich der Ausdruck für die freie Energie, die Ableitungen lauten:
\begin{align*}
f&=-J-\frac{1}{\beta}\log\klam\left[ {1+e^{-2\beta J}} \right]\\
\partial_\beta f&=\frac{1}{\beta^2}\log\klam\left[ {1+e^{-2\beta J}} \right]+\frac{2J}{\beta}\frac{1}{1+e^{2\beta J}}
\end{align*}

\end{document}